% !TEX root = ../main.tex
\chapteropening{%
Could we build a shortcut to the stars? General relativity turns that old
question into an engineering problem with a precise form: what geometry of
space and time would do the job, what would it take to produce that geometry,
and can nature supply what it takes?

This chapter is a tour of the problem with almost no mathematics. Every idea
in it returns later in equations, and most of them you will derive
yourself.}{%
\item what it means to engineer spacetime, and how that differs from
  engineering inside a geometry that nature has already fixed;
\item how hard it is to bend spacetime, and what sets that difficulty;
\item how a spacetime engineer reads Einstein's equation backwards;
\item the landmark designs of the field: traversable wormholes, warp drives
  and Krasnikov tubes;
\item why those designs are hard to build: the matter they need, the limits
  quantum physics places on it, and the threats to cause and effect.}
\chapter{What Is Spacetime Engineering?}\label{ch:what-is}

% ===========================================================================
\section{Clocks that disagree}\label{sec:gps}

\begin{thought}{The satellite's clock}
Two identical atomic clocks are set to the same time on the ground. One stays
there. The other is carried up to a navigation satellite, which circles the
Earth about \SI{20000}{\kilo\metre} overhead at nearly
\SI{4}{\kilo\metre\per\second}. Relativity offers two competing effects: a
moving clock runs slow, and a clock higher up in the Earth's gravity runs
fast. After a day, the two clocks are compared by radio. Which one is ahead?
Would the answer be the same for a clock on the space station, a few hundred
kilometres up?
\end{thought}

Every time a phone finds its position by satellite, it trusts clocks orbiting
far overhead. Each navigation satellite broadcasts the time on its onboard
atomic clock, and your receiver works out how far away the satellite is from
how long the signal took to arrive. Light travels about
\SI{30}{\centi\metre} in a nanosecond, so the satellite clocks must agree with
one another, and with clocks on the ground, to within a few tens of
nanoseconds if your position is to be good to a few metres.

At that precision, relativity cannot be ignored. Identical clocks in
different places do not tick at the same rate. A clock moving relative to you
runs slow, by an amount that grows with its speed. A clock higher up in a
gravitational field runs fast compared with a clock lower down: gravity
affects the rate of time itself. At everyday speeds and heights both effects
are minute, about a part in $10^{12}$ or less, and even for a satellite they
are only a few parts in $10^{10}$. A navigation clock, though, must keep time
to about one part in $10^{13}$.

\Cref{fig:clock-offsets} shows both effects for a clock on a circular orbit,
as a function of the orbit's size. The gain from height grows as the orbit
gets larger and levels off far from the Earth. The loss from motion is
largest for small orbits, because a low orbit must be flown fastest to stay
up. The net effect is the sum of the two, and it changes sign at an orbital
radius of one and a half Earth radii.

\begin{figure}[tb]
\centering
\includegraphics{ch01-clock-offsets}
\caption{How much a clock in a circular orbit gains per day over a clock on
the ground. Height makes the orbiting clock run fast (blue); motion makes it
run slow (orange). You will derive all three curves in \cref{ch:geometry}.
The Earth's rotation, which changes the numbers slightly, is ignored.}
\label{fig:clock-offsets}
\end{figure}

For the navigation satellites the net gain is about 38 microseconds per day.
Left uncorrected, it would spoil positions quickly: light travels
\SI{11}{\kilo\metre} in 38 microseconds, so the ranging error would grow by
about that much every day. The designers of the system corrected for it
before launch, by setting each satellite clock to tick slow by
$4.465$ parts in $10^{10}$ so that, once in orbit, it keeps pace with clocks
on the ground.

Navigation engineers work \emph{within} a geometry that nature has fixed. The
rates of clocks at different heights are set by the mass of the Earth, and the
engineers' job is to measure those rates and correct for them. Spacetime
engineering asks the reverse question. Suppose we could choose the geometry.
What pattern of clock rates, distances and light paths would we want? What
would such a spacetime do to travellers, to light and to ordinary matter? And
what would it take to produce it?

\begin{revisited}{The satellite's clock}
The satellite clock is ahead, by about 38 microseconds a day. At the
satellites' height, the gain from being higher (about 46 microseconds a day)
beats the loss from moving (about 7). The space station is barely higher than
the ground, about \SI{420}{\kilo\metre} up, but moves at
\SI{7.7}{\kilo\metre\per\second}; there the loss wins, and the station's clock
falls behind the ground by about 25 microseconds a day. Astronauts on the
station age slightly less than the rest of us.
\end{revisited}

% ===========================================================================
\section{How hard is it to bend spacetime?}\label{sec:stiff}

\begin{thought}{Denting time in the laboratory}
Popular pictures show gravity as a heavy ball denting a stretched rubber
sheet. Try making a dent yourself. Cast a ball one metre across from osmium,
the densest metal known: about twelve tonnes of it. Set one of today's best
atomic clocks against its surface, and an identical clock at the far end of
the laboratory, well away from it. Such clocks keep time to better than one
part in $10^{18}$. Will they detect the dent the ball makes in the rate of
time?
\end{thought}

Special relativity teaches that space and time are not separate. Observers in
relative motion disagree about lengths and durations, but they agree on the
speed of light and on a combination of time and distance called the
\emph{interval} between two events. Space and time together form a
four-dimensional \emph{spacetime}, and the interval plays the role in it that
distance plays in ordinary geometry.

General relativity takes one step further: the geometry of spacetime can be
curved, and gravity \emph{is} that curvature. An apple falls because it
follows the straightest possible path through a curved spacetime, and clocks
at different heights tick at different rates because the geometry of time
itself varies from place to place. John Wheeler summed up the theory in one
sentence: spacetime tells matter how to move, and matter tells spacetime how
to curve.

The rule for the second half of that sentence is \emph{Einstein's equation}.
In words, it says
\begin{equation}
  \begin{pmatrix}\text{curvature}\\ \text{of spacetime}\end{pmatrix}
  = \frac{8\pi G}{c^{4}}
  \begin{pmatrix}\text{energy, momentum,}\\ \text{pressure and stress}\\
  \text{of matter}\end{pmatrix}.
  \label{eq:einstein-words}
\end{equation}
The left side is built from the geometry. The right side describes
everything that has energy: matter, light and fields, and even the pressure
and tension inside a material. \Cref{ch:matter-geometry} turns both sides into
precise mathematical objects.

The constant in front, $8\pi G/c^{4}$, sets the scale of the whole subject.
Its size is fixed by
\begin{equation}
  \frac{c^{4}}{G} \approx \SI{1.2e44}{\newton},
  \label{eq:stiffness}
\end{equation}
a force, and an enormous one. Read \cref{eq:einstein-words} as a rule for
estimates: to curve spacetime so that its radius of curvature is about $L$,
matter must supply energy densities or stresses of order
$c^{4}/(8\pi G L^{2})$. For $L = \SI{1}{\metre}$ that is about
\SI{5e42}{\pascal}, more than thirty orders of magnitude beyond the breaking
strength of the strongest materials known. Spacetime is extraordinarily stiff.
That single number explains why the geometry around you is so nearly flat, and
it is the first reason spacetime engineering is hard.

A compact mass $M$ slows clocks near it by a fraction that, at a distance $r$
where the effect is small, is
\begin{equation}
  \frac{GM}{c^{2}r} = \frac{M}{r}\Big/\frac{c^{2}}{G},
  \qquad
  \frac{c^{2}}{G} \approx \SI{1.35e27}{\kilogram\per\metre}.
  \label{eq:clock-slowing}
\end{equation}
The ratio $c^{2}/G$ is the stiffness again, expressed as a mass per unit
length. A mass bends time appreciably only when the mass per unit distance
approaches that value.

\begin{example}[label=ex:earth-dent]{How far does the Earth bend time?}
Use \cref{eq:clock-slowing} to find how much slower a clock on the Earth's
surface runs than a clock far out in space. The Earth has mass
$M = \SI{5.97e24}{\kilogram}$ and radius $R = \SI{6.37e6}{\metre}$.

\solution The mass per unit distance is
$M/R = \SI{9.4e17}{\kilogram\per\metre}$. Dividing by $c^{2}/G$ gives
\begin{equation*}
  \frac{GM}{c^{2}R} = \frac{\SI{9.4e17}{}}{\SI{1.35e27}{}} = 7.0\times10^{-10}.
\end{equation*}
A clock on the ground loses 7 parts in $10^{10}$, about 60 microseconds a day,
compared with a clock far away.

\discussion The entire Earth produces less than a part-per-billion effect.
This is the same 60 microseconds a day at which the height curve of
\cref{fig:clock-offsets} levels off far from the Earth.
\end{example}

\begin{checkpoint}[label=chk:sun-ns]
Estimate the fractional slowing of clocks at the surface of the Sun
($M = \SI{2.0e30}{\kilogram}$, $R = \SI{7.0e8}{\metre}$), of a white dwarf of
0.6 solar masses and radius \SI{7000}{\kilo\metre}, and of a neutron star of
1.4 solar masses and radius \SI{12}{\kilo\metre}. For which of them is
\cref{eq:clock-slowing} a safe approximation?
\end{checkpoint}

\begin{revisited}{Denting time in the laboratory}
Not even close. The ball's mass per unit radius is
$\SI{1.2e4}{\kilogram}/\SI{0.5}{\metre} = \SI{2.4e4}{\kilogram\per\metre}$,
and dividing by $c^{2}/G$ in \cref{eq:clock-slowing} gives a slowing of about
$2\times10^{-23}$ at its surface. The clock across the laboratory is slowed by
only a few percent of that, so the difference between the two is still about
$2\times10^{-23}$, tens of thousands of times too small for the best clocks to
see. For the dent to reach one part in $10^{18}$, an osmium ball would have to
be about \SI{240}{\metre} across, some 160 million tonnes. The sheet is that
stiff.
\end{revisited}

% ===========================================================================
\section{Reading the equation backwards}\label{sec:backwards-intro}

\begin{thought}{Can relativity say no?}
A friend hands you a formula for a spacetime in which a ship crosses the
galaxy in an afternoon, and asks whether general relativity allows it.
Einstein's equation is the law that connects geometry to matter. Before you
read on, decide: is there any geometry that Einstein's equation, on its own,
forbids?
\end{thought}

Physicists usually read Einstein's equation from right to left. They describe
the matter first (a star, a galaxy, the contents of the universe) and then
solve for the geometry it produces. This is how the geometry around the Sun
was found, and the geometry of black holes and of the expanding universe. It
is hard work: the equation is a set of coupled, nonlinear partial differential
equations, and only highly symmetric problems can be solved by hand.

A spacetime engineer reads the equation from left to right. Start by deciding
what you want the spacetime to do: carry a traveller quickly between two
places, keep the traveller's clock in step with clocks at home, spare the
traveller from tidal stretching. Write down a geometry that does those things.
Then compute its curvature and read off from \cref{eq:einstein-words} the
matter that would have to be present to produce it. We call that matter the
\emph{demanded} matter of the geometry. In this direction the mathematics is
straightforward: computing curvature from a given geometry takes only
differentiation.

Michael Morris and Kip Thorne introduced this way of working in 1988, in a
paper written for students, to design wormholes a person could travel
through. They first listed the properties a traveller would need, then chose
a geometry with those properties, and only then asked what matter it would
require. The method is now called \emph{metric engineering}, because a
geometry is specified by its \emph{metric}, the object you will meet in
\cref{ch:geometry}.

The catch comes at the last step. Einstein's equation always returns an
answer, and nothing guarantees that the matter it names exists. The demand
might call for negative energy densities, or tensions far beyond any
material's strength, or both. The real work of spacetime engineering begins
there, and it runs as a loop (\cref{fig:design-loop}): test the demand
against everything known about matter, and revise the geometry until the
demand is one that nature might meet.

\begin{figure*}[tb]
\centering
\includegraphics{ch01-design-loop}
\caption{The design loop of spacetime engineering. A geometry is chosen for
what it does, and the matter it demands is computed and tested. When a test
fails, or when no known matter or field could supply the demand, the geometry
is revised. Most of this book teaches the steps inside this loop.}
\label{fig:design-loop}
\end{figure*}

The loop can also be entered from the other side. Some recent work starts
from matter with sensible properties, such as a shell of ordinary material,
solves for the geometry it produces, and then asks what that geometry can do
for a traveller. Both directions appear in this book. The geometry-first
direction is the more common, and it shows most clearly what a design costs.

\begin{checkpoint}[label=chk:division]
Team~A writes down a geometry that carries a ship to a nearby star in a week,
and computes the matter it demands. Team~B takes a moving shell of ordinary
steel and solves Einstein's equation for the geometry it produces. What is
each team sure of, and what must each still find out?
\end{checkpoint}

\begin{revisited}{Can relativity say no?}
Einstein's equation on its own forbids nothing. Read backwards, it assigns
every geometry the matter that would produce it, so it always has an answer.
The question \enquote{does relativity allow this?} therefore becomes
\enquote{does the matter this geometry demands exist, and can it be put in
place?} Your friend's afternoon crossing is possible only if nature can
supply its demand and put it in place, and that is where the physics bites.
\end{revisited}

% ===========================================================================
\section{Three ways to beat a light beam}\label{sec:gallery}

\begin{thought}{The race with a light beam}
A ship and a pulse of light leave the Earth together, bound for a star ten
light-years away. The light travels through ordinary empty space. The ship
reaches the star, turns round, and is home again before the pulse has even
arrived at the star. Yet relativity insists that nothing can overtake a light
signal travelling beside it. How can both be true? Try to invent a way
yourself before reading about the three that physicists have studied.
\end{thought}

\subsection{Tunnels through space}

A wormhole is a tunnel through space: two openings, called mouths, perhaps
light-years apart, joined by a short passage. \Cref{fig:wormhole} shows the
usual way of picturing one. Take a two-dimensional slice through space around
the wormhole and draw it as a curved surface. Far from the wormhole the
surface is flat. Near it, the surface curves into a funnel, narrows to a
smallest circle called the \emph{throat}, and opens out again into a second
region of space.

\begin{figure}[tb]
\centering
\includegraphics{ch01-wormhole}
\caption{An embedding diagram of a wormhole. A two-dimensional slice of space
is drawn as a surface. Circles around the wormhole shrink toward the throat
and grow again on the far side. A traveller (green) passes from region~A to
region~B through the throat.}
\label{fig:wormhole}
\end{figure}

Wormholes have appeared as solutions of general relativity since the 1930s,
but the early examples could not be crossed: they pinch off faster than
anything can pass through. Morris and Thorne asked what a wormhole would need
to be \emph{traversable}: no horizon at the throat, tidal forces a person
could survive, and a crossing that takes a reasonable time. They found
geometries that meet all these requirements, and then computed what those
geometries demand.

\begin{history}{a novel and a wormhole}
Morris and Thorne took up traversable wormholes after the astronomer Carl
Sagan sent Thorne the draft of his novel \emph{Contact} in 1985 and asked for
help making its physics as accurate as possible. Sagan revised the novel's
method of travel in the light of their preliminary results. Their paper,
subtitled \enquote{a tool for teaching general relativity}, launched the
modern study of engineered spacetimes.
\end{history}

The difficulty lies at the throat. A bundle of light rays heading into a
wormhole converges as it approaches the throat, and to come out the other side
it must spread apart again: the throat has to act on light as a
\emph{diverging} lens. What kind of matter could do that is the question of
the next topic.

\subsection{Moving space itself}

In 1994 Miguel Alcubierre took a different approach. Instead of joining two
distant places by a tunnel, his design moves a region of space itself. Space
contracts in front of a bubble and expands behind it
(\cref{fig:warp}a), and a region of flat spacetime inside the bubble is
carried along. A ship inside that region does not move through its
surroundings at all; it floats in free fall while the bubble carries it.

\begin{figure*}[tb]
\centering
\includegraphics{ch01-warp}
\caption{Alcubierre's warp bubble, in a plane containing the direction of
travel. (a)~The rate at which the volume of space grows (up, orange) or
shrinks (down, blue): space contracts ahead of the bubble and expands behind
it. (b)~The energy density measured by observers carried along by the flow of
space, who far from the bubble are at rest relative to the stars: negative
everywhere it is not zero, and concentrated in a ring around the bubble's
wall.}
\label{fig:warp}
\end{figure*}

Alcubierre's bubble has striking properties. Nothing inside it ever moves
faster than light relative to its own surroundings, yet the bubble as a whole
can move at any speed relative to distant stars. The ship feels no
acceleration, and its clock ticks at the same rate as clocks far away, so the
travellers age exactly as much during the trip as the people they left
behind. José Natário later showed that the expansion and contraction are not
essential: a warp drive can carry its passengers with no change of volume at
all. What \emph{is} essential is the demand. Observers carried along by the
flow of space find a negative energy density in the bubble's wall
(\cref{fig:warp}b), much as observers passing quickly through a wormhole's
throat would find one there.

\begin{checkpoint}[label=chk:warp-aging]
A crew makes a round trip to a star ten light-years away, once in a warp
bubble that moves at five times the speed of light, and once in a ship that
coasts at 99~percent of the speed of light, whose clocks run sevenfold slow.
For each trip, how much does the crew age, and how long do the people at home
wait?
\end{checkpoint}

\subsection{Building the road on the way out}

Sergei Krasnikov pointed out a problem with building and steering warp
drives, which the end of this chapter takes up. His alternative has the
traveller prepare the route on the way out. A ship travels to a distant star at less than the speed of light, and as
it goes it modifies the spacetime along its path, leaving behind a tube. The
modified geometry inside the tube lets the return trip arrive home shortly
after the departure. Allen Everett and Thomas Roman analysed these tubes
further and showed that two of them, suitably placed, would form a time
machine.

\subsection{What recent work has added}

In 2021 several authors proposed warp drives that appeared to need only
positive energy. Examining them closely, Jessica Santiago, Sebastian Schuster
and Matt Visser showed that the positive energy was measured by one special
family of observers only. Other observers passing through the same bubble
measure negative energy, and every physically reasonable warp drive violates
the conditions that ordinary matter satisfies. The episode taught the field a
lesson you will meet throughout this book: the demand of a geometry must be
checked from the point of view of \emph{every} observer. Current research
includes shells of ordinary matter that carry a flat interior at speeds below
that of light, computer simulations of warp spacetimes, and the search for
fields that could supply negative energy.

\begin{revisited}{The race with a light beam}
The light pulse never loses a race run side by side. In every design the ship
stays inside its local \emph{light cone}, the set of directions light can take
from where the ship is, so a light signal sent beside the ship would always
outrun it. What the designs change is the geometry that the ship and the pulse
travel through (\cref{fig:three-ways}). The wormhole shortens the ship's
route, so the pulse through ordinary space runs a longer race. The warp drive
moves the space the ship sits in: inside the bubble the light cones themselves
are tilted toward the ship's destination, and the ship rides inside them. The
Krasnikov tube tips the light cones along a route that the outbound trip has
prepared. There the ship even reaches the star after the pulse does, and yet,
by the clocks at home, it is back before the pulse has reached the star.
\end{revisited}

\begin{figure*}[tb]
\centering
\includegraphics{ch01-three-ways}
\caption{Three ways to reach a star and come home before a light pulse sent
at departure has reached the star, drawn as spacetime diagrams: distance runs
to the right, time runs up, and light travels at \ang{45}. Each wedge shows
the directions light can take from that event; the ship (green) always
travels inside the wedges. (a)~A wormhole shortens the route. (b)~A warp
bubble tilts the light cones inside it. (c)~A Krasnikov tube, laid down by the
outbound trip, tips the cones along the route.}
\label{fig:three-ways}
\end{figure*}

% ===========================================================================
\section{Matter that bends light the wrong way}\label{sec:obstacle}

\begin{thought}{A concave lens made of gravity}
Glass bends light either way: a convex lens focuses a beam, a concave lens
spreads it. Gravity bends light too, and the Sun's gravity focuses starlight
like a weak convex lens. A wormhole throat must spread the light that passes
through it, like a concave lens. Could you build a concave gravitational lens
out of ordinary matter, say from a ring or a hollow shell of very dense
material?
\end{thought}

Gravity bends light toward mass, and the Sun, like every ordinary body, acts
on light as a converging lens (\cref{ex:sun-lens}). Physicists summarize what
ordinary matter can do with a small set of inequalities called
\emph{energy conditions}. The simplest to state is the
\emph{weak energy condition}: every observer, however they move, measures an
energy density that is zero or positive. The most important for this book is
the \emph{null energy condition}, which is weaker still. It concerns light
passing through matter, and it guarantees that the matter's gravity can focus
a bundle of light rays but can never make a shrinking bundle start to grow
again before its rays cross (\cref{fig:focusing}a). Everything you have ever touched satisfies these
conditions, and so do light itself, electric and magnetic fields, and all
ordinary fluids.

\begin{example}[label=ex:sun-lens]{The Sun as a lens}
General relativity predicts that light passing a mass $M$ at distance $b$ is
bent toward it by the angle $4GM/(c^{2}b)$, twice what a Newtonian argument
gives. How strongly does the Sun bend light that grazes its edge, and where
does that light cross the axis on the far side? For the Sun,
$GM/c^{2} = \SI{1.48}{\kilo\metre}$ and $R = \SI{6.96e8}{\metre}$.

\solution For light grazing the edge, $b = R$ and the angle is
\begin{equation*}
  \frac{4 \times \SI{1.48e3}{\metre}}{\SI{6.96e8}{\metre}} = \SI{8.5e-6}{\radian},
\end{equation*}
or 1.75 seconds of arc. A ray bent inward by this small angle reaches the axis
after a distance $R$ divided by the angle, \SI{8.2e13}{\metre}, which is about
550 times the distance from the Earth to the Sun.

\discussion The deflection was first measured during the total solar eclipse
of 1919. Rays passing farther from the Sun are bent less and cross the axis
farther away, so the Sun focuses light along a line that starts some
550~astronomical units out, fourteen times as far as Pluto. The lens is weak,
and it is always a converging one.
\end{example}

\begin{figure}[tb]
\centering
\includegraphics{ch01-focusing}
\caption{(a)~Light passing through ordinary matter is focused: the bundle's
cross-section shrinks, and once it has started to shrink it keeps shrinking
until its rays cross.
(b)~Light passing through a wormhole throat follows the walls: the bundle
shrinks toward the throat and must grow again beyond it.}
\label{fig:focusing}
\end{figure}

The null energy condition is exactly what a shortcut has to violate. A
wormhole throat must make a shrinking bundle of light grow again
(\cref{fig:focusing}b), and a warp bubble's wall must do the same. Morris and
Thorne showed that a throat requires the material there to carry a tension
greater than its energy density, so that observers moving through the throat
close to the speed of light would measure a \emph{negative} energy density.
Matter with such properties is called \emph{exotic}, and you will derive this
requirement yourself in Part~II. In Chapter~5 you will see that it is no
accident. Under quite general
assumptions, any spacetime that lets travellers outrun light through an
otherwise ordinary region of space violates the null energy condition
somewhere.

\begin{keyidea}
Designing a shortcut through spacetime is easy on paper. Supplying the matter
it demands is the hard part: under very general conditions, a shortcut
demands matter that violates the null energy condition, and ordinary matter
never does.
\end{keyidea}

\begin{revisited}{A concave lens made of gravity}
No arrangement of ordinary matter will do. A glass lens spreads light by its
shape: glass slows light, and a concave lens is thicker at its edges than at
its centre. Gravity has no such freedom. The focusing power of matter is set
by the energy density and pressure of the matter that the light passes
through, and for ordinary matter it always focuses. A ring does not help.
Light passing through its middle is tugged one way and then the other, and
along its whole path the tugs cancel exactly, so it emerges undeflected, while
light passing outside the ring is bent inward. Light crossing a shell passes
through its walls, and the walls focus it. The tidal pull of nearby matter can
squeeze a bundle into an ellipse, but it never makes the bundle's area start
to grow. Making a shrinking bundle grow again before its rays cross takes
matter that violates the null energy condition.
\end{revisited}

% ===========================================================================
\section{Negative energy in the laboratory}\label{sec:casimir}

\begin{thought}{A stack of Casimir plates}
Quantum physics allows negative energy densities: between two closely spaced
metal plates, the energy density of the vacuum is lower than that of empty
space. Suppose you want to hold a wormhole's throat open with this negative
energy. Why not simply build an enormous stack of plates?
\end{thought}

In 1948 Hendrik Casimir predicted that two parallel, uncharged metal plates in
vacuum exert a force on each other. The plates restrict which electromagnetic vacuum
fluctuations fit between them, and the energy density between the plates is
lower than that of empty space; measured from empty space, it is negative.
For perfectly conducting plates a distance $d$ apart it is
\begin{equation}
  u = -\frac{\pi^{2}\hbar c}{720\,d^{4}},
  \label{eq:casimir}
\end{equation}
a result you will derive in Chapter~6. Steve Lamoreaux measured the Casimir
force in 1997, and it has since been measured many times. Quantum physics
therefore allows negative energy density.

It allows it only in limited amounts. Lawrence Ford and Thomas Roman showed
that quantum fields obey \emph{quantum inequalities}: the more negative the
energy density in a region, the shorter the time it can last and the smaller
the region it can occupy. Applied to a warp bubble, the inequalities force
the bubble's wall to be extraordinarily thin. Michael Pfenning and Ford
estimated that a bubble \SI{100}{\metre} in radius would then need about
\SI{6e62}{\kilogram} of negative energy for every unit of its speed in units
of $c$, vastly more than the mass of the observable universe. Clever
reshaping of the bubble reduces the requirement dramatically: Chris Van Den
Broeck brought it down to a few solar masses. That is still an astronomical
amount, and his design bends space on scales of about ten Planck lengths, far
below anything that has ever been probed.

\begin{checkpoint}[label=chk:casimir-force]
The energy between the plates is negative. Does that make the plates repel
each other? The energy between them, per unit area of plate, is $u\,d$: as the
plates move closer together, does it rise or fall, which way does the force
point, and how does its strength depend on $d$?
\end{checkpoint}

\begin{revisited}{A stack of Casimir plates}
A stack adds more gaps, but never deeper ones. A gap \SI{1}{\micro\metre} wide
holds about \SI{-4e-4}{\joule\per\cubic\metre}, however many plates surround
it, while a wormhole throat one metre in radius needs a tension exceeding its
energy density by roughly \SI{5e42}{\joule\per\cubic\metre}, some 46 orders of
magnitude more. Because the Casimir energy density grows as $1/d^{4}$,
reaching that value would take plates about \SI{3e-18}{\metre} apart, a few
hundred times smaller than a proton. This is the pattern the quantum
inequalities describe: large negative energy densities are confined to tiny
regions. The plates themselves, besides, carry the positive rest energy of the
metal, about \SI{2e20}{\joule\per\cubic\metre} for aluminium, which outweighs
the gap beside each one by more than twenty orders of magnitude.
\end{revisited}

% ===========================================================================
\section{Cause, effect and control}\label{sec:causality}

\begin{thought}{A message to the front}
You are aboard a warp bubble crossing space at twice the speed of light, and
you decide to stop. The bubble's front wall is only a hundred metres ahead of
you, so you send a radio command forward to shut it down. Does the command
arrive?
\end{thought}

Two further obstacles have nothing to do with the amount of energy. The first
is causality. Relativity allows observers in relative motion to disagree about
which of two distant events happened first, whenever the events are too far
apart for light to connect them. If a signal or a traveller could outrun
light across ordinary space, some observers would therefore see it arrive
before it left. With two faster-than-light links, suitably arranged, a
traveller could return before departing and meet an earlier self. Stephen
Hawking conjectured in 1992 that the laws of physics always prevent this,
a proposal he called the \emph{chronology protection conjecture}. It remains
a conjecture, and a responsible design must show that it creates no such
loops.

The second obstacle is control. A warp bubble is a flow of space. Inside it
the ship floats at rest in its own patch of space; across the bubble's wall
the flow changes to match the space of the distant stars; and light always
moves at the speed of light relative to the space around it. Krasnikov showed
what follows when the bubble outruns light: the crew is cut off from the front
of its own bubble, and can neither build the bubble nor steer it from inside.
Designs that carry travellers faster than light must therefore be laid out in
advance, or driven by something outside the travellers' reach.

\begin{checkpoint}[label=chk:krasnikov-control]
Robot ships leave the Earth at 90~percent of the speed of light for a star ten
light-years away, programmed to shape the wall of a warp bubble along the route
as they pass. Can the crew of a bubble that follows them steer it? How soon,
by the clocks at home, can the first bubble reach the star?
\end{checkpoint}

\begin{revisited}{A message to the front}
It never arrives, however close the wall is. Seen from the bubble, the space outside streams backward
past it at twice the speed of light, and inside the front wall the stream
slows smoothly to zero. Your radio signal moves forward at the speed of light
relative to the space it is in, so somewhere in the wall, where the backward
stream reaches the speed of light, the signal stalls and can advance no
further. The same thing happens to light trying to leave a black hole, and in
\cref{ch:geometry} you will see that this place is a horizon.
\end{revisited}

% ===========================================================================
\section{Where the subject stands}\label{sec:today}

Nobody knows how to build any of the spacetimes in this chapter, and some of
them may be forbidden outright. Spacetime engineering today is a theoretical
discipline, and it earns its place in three ways.

It teaches general relativity. Designing a spacetime forces you to understand
how every piece of its geometry affects clocks, rulers, light and matter. That
was Morris and Thorne's reason for writing their paper, and it is one reason
for this book.

It maps the boundary of the possible. Some of the obstacles above are theorems
with precise assumptions, some are estimates that depend on how a design is
drawn, and some are conjectures. Knowing which is which tells you what would
have to change for a design to become possible, and much of this book is
about learning to tell them apart.

It develops methods that carry over to other problems: how to design under
severe physical constraints, how to check a long calculation, and how to judge
how far a result can be trusted.

The rest of Part~I gives you the tools. \Cref{ch:geometry} develops the
geometry of spacetime: intervals, clocks, the metric, free fall and
curvature. \Cref{ch:matter-geometry} describes matter by its energy, momentum
and stress, states Einstein's equation, and shows how to compute what a
geometry demands. The chapters after them split spacetime into space and time
the way an engineer describes a design, state the energy conditions
precisely, explain what quantum physics allows, and treat causality. Part~II
studies the classic designs in detail. Part~III is about design itself: how
the rate of clocks, the flow of space and the shape of space each affect what
a spacetime demands and what it does to the people inside it. Part~IV asks
what matter and fields could supply a demand, and Part~V takes up dynamics,
the checking of calculations, and the open problems of the field.

This chapter used SI units. From \cref{ch:geometry} on we measure time and
distance in the same units, which sets $c = 1$, and from
\cref{ch:matter-geometry} on we also set $G = 1$. Whenever a number matters,
we convert back to everyday units.

% ===========================================================================
\begin{chaptersummary}
\item Relativity makes identical clocks tick at different rates in different
  places. Satellite navigation works only because its designers correct for
  both the speed and the height of the satellites.
\item Spacetime has a geometry that matter bends; gravity is that bending.
  Einstein's equation links the curvature of spacetime to the energy,
  momentum, pressure and stress of matter.
\item The constant $c^{4}/G \approx \SI{1.2e44}{\newton}$ makes spacetime
  extremely stiff: curving it with a radius of curvature $L$ takes energy
  densities or stresses of order $c^{4}/(8\pi G L^{2})$. A mass $M$ slows
  clocks at distance $r$ by about
  $GM/(c^{2}r)$, which is small unless $M/r$ approaches
  $c^{2}/G \approx \SI{1.35e27}{\kilogram\per\metre}$.
\item Spacetime engineering reads Einstein's equation backwards: choose a
  geometry for what it does, then compute the matter it demands. Every
  geometry has a demand; whether nature can meet it is the real question.
\item The classic designs are traversable wormholes (a shorter route), warp
  drives (space that carries the ship) and Krasnikov tubes (a route prepared
  on the way out). In each, a traveller always moves inside the local light
  cone.
\item All of them demand matter that violates the null energy condition,
  which ordinary matter satisfies. Quantum physics allows negative energy
  density, as the Casimir effect shows, but quantum inequalities limit how
  much, for how long and over how large a region.
\item Faster-than-light designs also face causality (the risk of time
  machines) and control (a bubble moving faster than light cannot be steered
  from inside).
\end{chaptersummary}

\begin{checkanswers}
\item[\ref{chk:sun-ns}] Sun: $GM/(c^{2}R) \approx 2\times10^{-6}$. White dwarf:
  about $1.3\times10^{-4}$. Both are small, so the estimate is safe. Neutron
  star: about $0.17$, not small, so the estimate is rough; the exact result of
  \cref{ch:geometry} is a clock rate of 0.81 of the rate far away. The step
  from star to white dwarf to neutron star is a step in mass per unit radius.
\item[\ref{chk:division}] Team~A is sure its geometry does the job, and it can
  always compute the demand, since that takes only differentiation; it must
  still find out whether matter with that demand can exist and be put in
  place. Team~B is sure its matter exists. It must solve coupled nonlinear
  equations, which is possible by hand only with a great deal of symmetry and
  needs conditions far from the shell as well, and it may find that its
  geometry does nothing useful for a traveller.
\item[\ref{chk:warp-aging}] Warp bubble: the round trip covers 20 light-years at
  five times the speed of light, so it takes 4 years by the clocks at home.
  The crew floats at rest in its own patch of space, and its clocks keep pace
  with those at home, so the crew also ages 4 years. Coasting ship: the trip
  takes $20/0.99 = 20.2$ years at home, while the crew, moving through the
  space around it, ages $20.2/7.09 = 2.85$ years. The fast ship shortens the
  trip only for its crew, who come home to a world 20 years older; the warp
  bubble shortens it for everyone.
\item[\ref{chk:casimir-force}] No: they attract.
  $u\,d = -\pi^{2}\hbar c/(720\,d^{3})$ grows more negative as $d$ shrinks, so
  the energy falls as the plates approach, and the plates move the way that
  lowers it. The force per unit area is the rate at which the energy falls
  with distance, $\pi^{2}\hbar c/(240\,d^{4})$: about \SI{1.3}{\milli\pascal}
  for plates \SI{1}{\micro\metre} apart, rising to a full atmosphere at about
  \SI{11}{\nano\metre}. Laboratory measurements of the attraction agree with
  this prediction.
\item[\ref{chk:krasnikov-control}] No. The robots follow a program fixed
  before they left, and nothing the crew does can reach the parts of the route
  ahead of the bubble. The bubble also cannot outrun the robots that prepare
  its path, so the first trip reaches the star no sooner than they do,
  $10/0.9 = 11.1$ years after they left: later than a light signal would.
  Only later trips gain, once the route is prepared. This is Krasnikov's
  point, and the idea behind his tube.
\end{checkanswers}

\begin{problems}
\item \diff{1} \textbf{Two clocks.} A clock sits at the top of a tower of height
  $h = \SI{100}{\metre}$. Near the Earth's surface, a clock higher by $h$ runs
  faster by the fraction $gh/c^{2}$, where
  $g = \SI{9.8}{\metre\per\second\squared}$. How many nanoseconds does the
  top clock gain over a clock at the base in one year?
\item \diff{1} \textbf{The stiffness of spacetime.} Compute $c^{4}/G$ in newtons
  from $c = \SI{2.998e8}{\metre\per\second}$ and
  $G = \SI{6.674e-11}{\newton\metre\squared\per\kilogram\squared}$. Then
  estimate the stress, in pascals, needed to curve spacetime appreciably over a
  distance of \SI{1}{\kilo\metre}, and compare it with the tensile strength of
  steel, about \SI{e9}{\pascal}.
\item \diff{1} \textbf{A season in orbit.} An astronaut spends six months on the
  space station. Using \cref{fig:clock-offsets}, estimate how much less the
  astronaut ages than someone who stayed on the ground.
\item \diff{2} \textbf{A day of navigation error.} Using the numbers in this
  chapter, find how long it would take an uncorrected navigation-satellite
  clock to drift by the \SI{30}{\nano\second} that corresponds to about
  \SI{10}{\metre} of position error.
\item \diff{2} \textbf{How much is too much?} Express the Pfenning--Ford estimate
  of \SI{6e62}{\kilogram} in solar masses, taking
  $\Msun = \SI{2.0e30}{\kilogram}$. Van Den Broeck's design needs a few solar
  masses. By what factor did reshaping the bubble reduce the requirement?
\item \diff{2} \textbf{Casimir numbers.} Evaluate \cref{eq:casimir} for plates
  \SI{100}{\nano\metre} apart. How thick a layer of this negative energy
  density, spread over the surface of a sphere of radius \SI{1}{\metre}, would
  contain a total negative energy equal to the rest energy of one gram of
  matter?
\item \diff{3} \textbf{Which assumptions?} The quantum inequalities rest on
  assumptions: the kind of quantum field, the flatness of spacetime over the
  sampling time, and the way the energy is averaged. For each, describe how a
  designer might try to take advantage of it, and what you would check to see
  whether the attempt succeeds. Return to your answer after Chapter~6.
\end{problems}

\begin{furtherreading}
\item[\textcite{Ashby2003}] The complete relativistic treatment of satellite
  navigation, including the clock corrections of the first topic.
\item[\textcite{Einstein1916}] Einstein's own presentation of general
  relativity, still readable today.
\item[\textcite{Dyson1920}] The eclipse expedition of 1919 that measured the
  Sun's bending of starlight.
\item[\textcite{Brewer2019}] An atomic clock that keeps time to better than one
  part in $10^{18}$, the precision assumed in the thought experiment on
  denting time.
\item[\textcite{Morris1988a}] The paper that started the field, written for
  students and still one of the best introductions to engineered spacetimes.
\item[\textcite{McMonigal2012,Schuster2023}] Two treatments of metric
  engineering as a method.
\item[\textcite{Alcubierre1994}] The five-page paper that introduced the warp
  drive, readable after \cref{ch:matter-geometry}.
\item[\textcite{Natario2002}] Warp drives without expansion.
\item[\textcite{Krasnikov1998,Everett1997}] Krasnikov's tube, and its
  analysis, including the time machine that two tubes would form.
\item[\textcite{Santiago2022}] Why every physically reasonable warp drive
  violates the energy conditions.
\item[\textcite{Olum1998}] Why travel faster than light through ordinary
  space requires negative energy.
\item[\textcite{Casimir1948,Lamoreaux1997}] The prediction of the Casimir
  effect and its first precise measurement.
\item[\textcite{Ford1995,Ford1997,Pfenning1997}] Quantum inequalities, and what
  they imply for wormholes and warp bubbles.
\item[\textcite{VanDenBroeck1999}] How reshaping a warp bubble reduces its
  energy requirement.
\item[\textcite{Hawking1992}] The chronology protection conjecture.
\item[\textcite{Everett2011}] A book-length, non-technical account of warp
  drives, wormholes and time travel by two researchers who helped shape the
  field.
\item[\textcite{Visser1995b}] The standard monograph on wormholes.
\item[\textcite{Kontou2020}] A modern review of energy conditions in classical
  and quantum physics, useful throughout this book.
\end{furtherreading}
